\section{Discussion}
\label{sec:discussion}
In this benchmark, which method appears best depends strongly on the protocol. This section discusses which elements of the protocol drive the conclusions (Section~\ref{sec:disc-drivers}), what the results imply for the design of two-phase hybrids (Section~\ref{sec:disc-hybrids}) and for layout practice (Section~\ref{sec:disc-practice}), and condenses them into a checklist for benchmarking (Section~\ref{sec:disc-recs}); Section~\ref{sec:limits} states the threats to validity.

\subsection{What Drives the Conclusions of Metaheuristic Comparisons}
\label{sec:disc-drivers}
\emph{Baseline configuration.} The largest single effect is the setting of the PSO baseline. An advantage over PSO is informative only if the PSO baseline is convergent: the old setting $w=0.7$, $c_1=c_2=2$ is order-1 stable but lies beyond the order-2 boundary~\eqref{eq:poli}, which marks approximately where this PSO starts to fail, and places PSO at position \NBaseOldPSOPos{} of Table~\ref{tab:baseline} instead of \NBaseNewPSOPos{} (Section~\ref{sec:psosetting}). Run without velocity clamping and with clipped coordinates keeping their velocity, its swarm contracts much less, far fewer of its late evaluations are feasible, and it returns a feasible layout in only \NOldPSOFeas\% of the runs. The sweep of $c_1=c_2=c$ connects this to Proposition~\ref{prop:pso-stability} without establishing a cause: at least \NSFeasBelow\% of the runs are feasible up to $c=\NSCBelow$, and feasibility starts to decline in the step that contains $c^\ast=\NSBoundC$ and then falls progressively (\NSFeasCTwoZero\% at $c=2$). Clamping the velocities largely restores feasibility (\NSFeasVMax\%), whereas zeroing clipped components does not (\NSFeasVZero\%), so the old setting fails mainly through large steps rather than through the bound handling. The stability region, derived under stagnation and without bounds, is therefore a cheap screen for implausible settings, not a guarantee of good behavior; the constriction setting~\eqref{eq:constriction} passes it and was fixed a priori from the literature, not tuned on our cases.
% CHECK-FINAL [C27]: Table baseline: old setting PSO position 5, constriction position 1; CHECK-CSWEEP [S04]/[S05]: feasibility loss begins in the step containing c* and accelerates ("marks approximately where ... starts to fail"; no causal claim, no jump)
% CHECK-THEORY [T01]/[T04]: order-1 stable (4 < 6.8), order-2 violated (4 > 3.50); CHECK-DIAG [D03]: "contracts much less"; [D05]: far fewer feasible particle evaluations in iterations 101-200; CHECK-FINAL [C12] + generated \NOldPSOFeas
% CHECK-CSWEEP [S02]: >= \NSFeasBelow % feasible for c <= \NSCBelow; [S03]-[S05]: decline starts in the step containing c*, progressive; [S09]: zeroing clipped velocity no rescue; [S10]/[S11]: vmax clamping \NSFeasVMax % ("largely"), still worse than constriction -> "mainly through large steps"

\emph{Controls.} An advantage of a hybrid's first phase is informative only if it beats random sampling in the farm at equal cost. Square sampling is too weak a control: \NRSRunsNoFeasSample{} RS-VNS runs draw no feasible sample, mainly because of the spacing constraint, so a salp-swarm phase appears to help against it. Against disc sampling (RSD-VNS, third of the \NAblNVar{} component-analysis variants), only the PSO phase gains: a gain of the SSA phase larger than the margin is ruled out across seeds and cases, LX-SSA as published is worse, and the small SSA gain over RS-VNS loses its significance with 1\textdegree{} bins (Section~\ref{sec:ablation}). As the second phase is essentially one compass-search descent, this measures the value of a start for a local search; without an in-site control of equal cost, the contribution of a first phase remains unknown~\cite{Feng2015}.
% CHECK-DIAG [D17] + CHECK-FINAL [C57]/[C60]: square control weak (spacing binding); [C58]: RSD-VNS ranks third; [C53]/[C54]: no SSA gain over disc sampling; [C55]: LX-SSA-VNS worse; [C56]: PSO-VNS beats RSD-VNS ("only the PSO phase gains"); CHECK-EXTRA [X53]: SSA gain > margin ruled out (seed, case); CHECK-DIR [F17]: SSA-VNS vs RS-VNS n.s. at 1 deg; CHECK-DIAG [D08]: essentially one descent

\emph{Level of equivalence.} Non-significance is not equivalence, and equivalence depends on the level and the evaluation: PSO-VNS and PSO are equivalent at $\pm\NEqMargin$~pp across seeds and cases but not across farm clusters, nor under a 1\textdegree{} rose, where PSO-VNS has a small significant advantage (Sections~\ref{sec:overall} and~\ref{sec:robustness}). The three levels answer different questions: the average over this fixed benchmark, over similar cases, or over new farms. The minimal margin grows accordingly (\NXEqSeedPSOVNSvsPSOMin, \NXEqCasePSOVNSvsPSOMin{} and \NXEqClustPSOVNSvsPSOMin~pp), and the average rests on the small layouts: the two methods are equivalent for $N<10$ (\NXEqSmallN{} cases) but not for $N\ge10$, and \NXHetBeyond{} of the \NXHetN{} cases differ by more than the margin, in both directions. An equivalence statement is thus a statement about a level, a margin and a benchmark composition; reporting the minimal margins lets readers apply their own threshold.
% CHECK-EXTRA [X39]/[X41]: equivalent at seed and case level, not at cluster level (CR2), minimal margins seed < case < m < cluster; CHECK-DIR [F10]: 1 deg -0.043 pp, 90% CI below 0, not equivalent; CHECK-EXTRA [X44]: N < 10 equivalent (\NXEqSmallN cases), N >= 10 not; [X43]: \NXHetBeyond of \NXHetN beyond +-m, both directions

\emph{Discretization of the objective.} Part of the ranking is a property of the 15\textdegree{} rose. Re-evaluated with 1\textdegree{} sub-bins, the mean wake loss of every method rises (PSO-VNS: \NFLossPSOVNSFifteen\% to \NFLossPSOVNSOne\%), that of PSO and the VNS-based methods most, because optimized layouts exploit the gaps between the coarse bins (Section~\ref{sec:robustness}). PSO-VNS and PSO keep the first two places, but the order beyond them changes (MS-SLSQP moves from sixth to third), PSO's narrow lead under the Gaussian wake with 15\textdegree{} bins disappears, and the equivalence verdict for PSO-VNS versus PSO turns into a small significant difference. From feasible starts, a single local search (VNS alone) ranks first on the six largest cases (Section~\ref{sec:feasinit}), and at 120,030 evaluations the leading methods lie within \NBudCloseGapOneTwentyK~pp (Section~\ref{sec:budget}). A ranking is therefore the ranking of a protocol: baseline settings, control, level of inference, evaluation resolution, initialization and budget.
% CHECK-DIR [F03]: loss rise of every method (PSO-VNS 1.322 -> 2.335); [F05]: PSO and VNS-based methods rise most; [F06]: PSO-VNS first, PSO the only one not significantly different; [F07]: MS-SLSQP 6th -> 3rd; [F09]: Gaussian: PSO first at 15 deg only; [F10]: not equivalent at 1 deg
% CHECK-FINAL [C28]: feasible starts, VNS alone first; [C37]/[C45]: 120,030 evaluations, leading methods within \NBudCloseGapOneTwentyK pp of PSO-VNS

\subsection{Implications for the Design of Two-Phase Swarm--Local-Search Hybrids}
\label{sec:disc-hybrids}
\emph{The swarm phase as a feasibility-reaching start.} In a relay hybrid, the first phase has to hand the local search a good start at low cost; it need not be a good optimizer. In our runs, the value of the PSO phase lies largely in reaching feasibility. At the switch, its global best is feasible in \NSwitchFeasPSOVNS\% of the runs, against \NSwitchFeasSSAVNS\% for SSA, \NSwitchFeasRSDVNS\% for disc sampling and \NSwitchFeasRSVNS\% for square sampling (Section~\ref{sec:ablation}); when feasible, the best disc sample has a wake loss of \NSwitchLossRSDVNS\% against \NSwitchLossPSOVNS\% for the PSO phase, averaged over different subsets of runs. From feasible starts, PSO never improves on the best initial layout, because its moves change all coordinates at once and make turbines collide, and VNS alone ranks first, with PSO-VNS second (Section~\ref{sec:feasinit}). A swarm phase is thus most useful when random initial layouts are infeasible; when a feasible start is cheap (our packing initialization needs no objective evaluations), the budget is better spent on the local search. Whether PSO also explores better cannot be separated from its faster feasibility with our data.
% CHECK-FINAL [C16]: PSO switch feasibility higher and switch loss lower than SSA / LX-SSA; [C57]: RSD-VNS feasible at the switch more often than RS-VNS; CHECK-DIAG [D18]: PSO > SSA > RS at the switch; switch macros from ablation.phase2_loss_reduction_pct (generated; the PSO-VNS vs RSD-VNS loss comparison is descriptive, different run subsets)
% CHECK-FINAL [C28]/[C38]: feasible starts, VNS alone first, PSO-VNS second, PSO and DE never move; CHECK-DIAG [D13]/[D14]: PSO never improves on the best initial layout, first-move collisions; CHECK-THEORY [T11]

\emph{The local-search phase.} At 6,030 evaluations, Phase~2 is essentially one compass-search descent: the first descent gives \NDDescentSharePSOVNS\% of the Phase-2 gain, it is unfinished in \NDDescentUnfinishedLargePSOVNS{} of \NDRunsLargePSOVNS{} runs with $N\ge12$, and the shake-and-search cycles take \NDCycleEvalsPctPSOVNS\% of the Phase-2 evaluations for \NDCycleSharePSOVNS\% of the gain (Section~\ref{sec:ablation}). Our conclusions about ``VNS'' are therefore conclusions about this local search. For the design of such hybrids at small budgets, the cost of the descent, $2n$ evaluations per sweep here, matters more than the neighborhood structure, which is exercised only for small $N$; cheaper descents (first improvement, single-turbine moves) are the obvious candidates, but we did not test them. The descent also repairs infeasible switch points at no extra cost: in the densest cases (500~m, $N=10$), the PSO global best is feasible at the switch in only \NDenseSwitchFeas\% of the runs, but all runs are feasible after Phase~2.
% CHECK-DIAG [D08]: first descent \NDDescentSharePSOVNS % of the Phase-2 gain; [D09]: N >= 12 unfinished \NDDescentUnfinishedLargePSOVNS of \NDRunsLargePSOVNS; [D10]: accepted cycles only in k = 1; [D11]: infeasible switch points repaired by the first descent; [D20]: cycle shares; CHECK-FINAL [C24]: dense_500_10 feasible at switch < 100, final 100

\emph{Budget split.} The split $\omega$ is the only parameter PSO-VNS adds to PSO and VNS. Of the splits tested on \NSplitCases{} cases, \NSplitBestSetting{} has the lowest mean wake loss and average rank, $\omega=0.9$ does not beat it, and plain PSO ($\omega=1$) has a higher mean loss (Section~\ref{sec:split}). This is consistent with the feasibility view: the swarm needs enough evaluations to reach the feasible region, and the local search enough to complete much of its first descent. The preset $\omega=0.5$ was not tuned; a larger split is a candidate to confirm on the target problem.
% CHECK-FINAL [C47]: 0.75 best (lowest mean loss and rank), omega = 1 higher mean loss than 0.75; [C59]: 0.9 does not beat 0.75; CHECK-EXTRA [X36]/[X37]

\emph{Where PSO-VNS helps and where it does not.} At 6,030 evaluations from random starts, PSO-VNS helps in three places. It reaches feasibility more reliably on the Horns Rev~1 block (\NHRFeasPSOVNS{} vs.\ \NHRFeasPSO{} runs, McNemar $p=\NXHRMcNemarP$), where the descent repairs residual infeasibility; on the benchmark, PSO is feasible in every run as well. It ranks best on the six largest cases at all three budgets (exploratory). And in the post hoc subgroup of the larger Data Set~II layouts ($N\ge10$), it has the lower mean wake loss in \NCmPSOVNSvsPSOdsIILargeWins{} of \NCmPSOVNSvsPSOdsIILargeN{} cases ($p_{\rm Holm}\le\NXHolmPSOVNSvsPSOdsIILargeP$; \NXEnPSOVNSvsPSOdsIILargePct\% of AEP), an advantage that remains with 1\textdegree{} bins. It does not help from feasible starts, in the densest Data Set~I layouts (all \NXHetDsINPSO{} Data Set~I cases beyond the margin favor PSO), on IEA37, where the best published layouts are better, or, averaged over the benchmark, by more than a few hundredths of a percentage point. PSO-VNS is thus a consistent reference method for this benchmark, not a superior optimizer.
% CHECK-EXTRA [X52]: Horns Rev 30/30 vs 19/30, McNemar p 9.8e-4 (one site, random starts); CHECK-FINAL [C08]/[C20]; generated \NFeasPSO = 100 (PSO feasible in every benchmark run; R5 O8); [C29]: six largest cases best at all budgets (exploratory)
% CHECK-FINAL [C32] + CHECK-EXTRA [X19]/[X27]: DS II N >= 10 \NCmPSOVNSvsPSOdsIILargeWins of \NCmPSOVNSvsPSOdsIILargeN, all-family Holm; CHECK-DIR [F14]: survives 1 deg; CHECK-EXTRA [X43]: all \NXHetDsINPSO DS I cases beyond m favour PSO; [X25]: mean gain < 0.05 % AEP; CHECK-FINAL [C28]/[C40]

\subsection{Implications for Wind Farm Layout Practice}
\label{sec:disc-practice}
\emph{Discretized roses are exploited.} The direction discretization is the most important caveat for practice. Layouts optimized under a 15\textdegree{} (benchmark) or 5\textdegree{} (Horns Rev~1) rose place turbines in the gaps between the bins: with 1\textdegree{} bins the benchmark wake loss rises by \NFLossRiseMin--\NFLossRiseMax~pp, more than any difference between the leading methods, and none of the \NFHRNFeas{} feasible optimized Horns Rev~1 layouts exceeds the installed block (\NFHRAboveFiveTotal{} with 5\textdegree{} bins), although the $4D$ spacing enlarges their feasible set. Energy gains of layouts optimized under a coarse rose should therefore be verified at a fine resolution before they are attributed to an optimizer, and the resolution should be refined during optimization when the cost per evaluation allows it.
% CHECK-DIR [F03]: every method's loss rises by \NFLossRiseMin to \NFLossRiseMax pp (about \NFLossRiseMarginRatio x the margin for PSO-VNS); [F21]: 0 of 901 feasible runs above the installed block at 1 deg (16 at 5 deg); CHECK-FINAL [C62]: no \NHR...Above macro used

\emph{Horns Rev~1 and the wake model.} The Horns Rev~1 block tests optimizers on a measured wind climate, not the installed design: it is evaluated without the other 64 turbines, with $4D$ instead of the installed $7D$ spacing, and without cables, foundations, loads or navigation constraints (Section~\ref{sec:hornsrev-site}). Our model takes $C_T$ at the free-stream speed, which gives a wake loss \NFPyWakeLossDiffPP~pp relative to PyWake's NOJ model at identical bins; with $C_T$ at the local speed, it reproduces PyWake's hub-center NOJ exactly. Design studies should evaluate $C_T$ at the local speed, and the benchmark's constant $C_T$ overstates the deficits above rated speed (Section~\ref{sec:model}). For Data Set~I, the benchmark's wake-free capacity factor is 62\%, and a cubic power curve lowers the expected power by \NCubicOverI\% (Data Set~I) and \NCubicOverII\% (Data Set~II; Section~\ref{sec:robustness}): the benchmark's energy values are not site estimates.
% CHECK-DIR [F32]/[F34]: PyWake 661.39 vs ours 666.75 GWh/yr, wake loss -0.72 pp, cause C_T at free-stream speed; local C_T reproduces PyWake's hub-centre NOJ; CHECK-FINAL [C21]: cubic curve (generated \NCubicOverI / \NCubicOverII); 62 % capacity factor verified in 03_model (14045.74/15/1500)

\emph{IEA37 and the gap to published layouts.} None of our methods reaches the best feasible published IEA37 layouts: at 30,030 evaluations the best PSO-VNS layouts are \NFIEAGapSixteenThirtyKProj\% (16 turbines) and \NFIEAGapThirtySixThirtyKProj\% (36 turbines) below them in AEP, and with 36 turbines the gradient-based MS-SLSQP gives our best layout (Section~\ref{sec:iea37}). The published layouts come from optimizers chosen by the participants, including gradient-based ones, with budgets that were not reported comparably~\cite{Baker2019,Thomas2023}. The metaheuristics compared here, at the budgets used, are therefore reference points for comparing optimizers, not design tools. The optimized benchmark layouts also often lie on the $4D$ constraint (only \NSpFiveLarge\% of those with $N=11$--15 satisfy $5D$) and do not transfer to larger spacings without re-optimization (Section~\ref{sec:spacing}).
% CHECK-DIR [F44]: projected IEA37 gaps -1.93 / -6.23 % (30,030); [F45] + CHECK-FINAL [C48]: 36 turbines, best of our layouts from MS-SLSQP; [C40]: best published layouts better; [C23]: spacing share (\NSpFiveLarge % of N = 11-15 satisfy 5D, < 10)

\subsection{Recommendations for Benchmarking}
\label{sec:disc-recs}
Table~\ref{tab:recommendations} condenses the protocol into a checklist for comparisons of layout optimizers~\cite{BartzBeielstein2020,LaTorre2021}; each item states why it matters in this benchmark and where the paper implements it.
\begin{table}[!htbp]
\centering
\caption{Checklist for benchmarking metaheuristics on the WFLOP: item, rationale from this study, and what we did (where in the paper).}
\label{tab:recommendations}
\footnotesize
\renewcommand{\arraystretch}{1.15}
\begin{tabular}{@{}>{\raggedright\arraybackslash}p{0.03\textwidth}>{\raggedright\arraybackslash}p{0.27\textwidth}>{\raggedright\arraybackslash}p{0.34\textwidth}>{\raggedright\arraybackslash}p{0.30\textwidth}@{}}
\toprule
& Item & Rationale & What we did (where) \\
\midrule
1 & Use convergent baseline settings (for PSO, \eqref{eq:poli} or \eqref{eq:constriction}); state bound and velocity handling; verify with instrumented runs & The old setting reverses the order of PSO and the salp-swarm hybrids; its failure depends on the velocity handling & Stability check, $c$-sweep, bound variants, instrumented reruns (Sections~\ref{sec:pso}, \ref{sec:psosetting}; Table~\ref{tab:baseline}) \\
2 & Give all methods the same budget, counting the initial population, gradient calls and repeated evaluations & Iteration budgets favor methods with cheap iterations; LX-SSA costs $2N_p$ per iteration, MS-SLSQP pays for gradients & $B$ fixed in evaluations for all methods (Sections~\ref{sec:template}, \ref{sec:setup}) \\
3 & Pair runs by seed with a common initial population & Pairing removes the differences between initial populations from all contrasts & Seeds 1--30 shared by all methods (Section~\ref{sec:setup}) \\
4 & Rank feasibility first and report feasibility rates & Penalized methods return infeasible layouts; means over feasible runs alone favor unreliable methods & Qualification threshold, feasibility-first ranks (Section~\ref{sec:setup}; Table~\ref{tab:friedman68}) \\
5 & Use paired nonparametric tests with multiplicity correction, on case means and runs & Mean-rank post hoc tests depend on the pool; many contrasts inflate errors & Friedman, Holm-corrected Wilcoxon tests, W/T/L (Sections~\ref{sec:setup}, \ref{sec:overall}) \\
6 & Report the robustness of inference to thresholds and case dependence & Cases share farms and data sets; one threshold can decide a verdict & Six thresholds, \NXClusters{} clusters, leave-one-out, one Holm family (Section~\ref{sec:setup}) \\
7 & Test equivalence at a stated level and margin; report minimal margins & Non-significance is not equivalence; verdicts change with the level & TOST at seed, case and cluster level, Bayesian test (Sections~\ref{sec:setup}, \ref{sec:overall}) \\
8 & In ablations of hybrids, include the local method alone and an in-site random-sampling control of equal cost & A weak (square) control makes an uninformative first phase look useful & VNS, RS-VNS and RSD-VNS controls (Section~\ref{sec:ablation}; Table~\ref{tab:ablation}) \\
9 & Vary budget and initialization before generalizing a ranking & Rankings at a small budget and from random starts need not hold otherwise & Three budgets, feasible starts (Section~\ref{sec:beyond}; Table~\ref{tab:feasbudget}) \\
10 & Re-evaluate final layouts at a finer direction resolution and under other models; publish code and per-run data & Optimized layouts exploit coarse bins; part of the ranking is a property of the discretization & 1\textdegree{} bins, Gaussian wake, cubic curve, PyWake, IEA37; public code and data (Sections~\ref{sec:beyond}, \ref{sec:robustness}) \\
\bottomrule
\end{tabular}
\end{table}
% CHECK-DIAG [D01]/[D02]: instrumented runs reproduce the stored runs; CHECK-CSWEEP [S09]/[S10]: bound/velocity handling changes the old setting's feasibility (item 1); CHECK-FINAL [C27] (item 1); CHECK-EXTRA [X02]-[X09], [X11]/[X14], [X18] (item 6); [X41]: minimal margin grows with the level (item 7); CHECK-FINAL [C57]/[C60] (item 8); [C28]/[C37] (item 9); CHECK-DIR [F03]/[F07] (item 10)

\section{Threats to Validity}
\label{sec:limits}
We group the limitations of the study by the usual four types of validity threats.

\emph{Internal validity} (are the observed differences due to the methods and components?). The comparisons of PSO-VNS with PSO and of each Phase~1 with the controls are clean by construction: PSO-VNS repeats PSO exactly up to the switch, and the controls differ from the hybrids only in Phase~1 (3,015 versus 3,030 evaluations). Our reused SSA, LX-SSA, PSO and DE code reproduces the result distributions of the original code, and the instrumented copies reproduce the stored runs (Section~\ref{sec:validation}). The configuration of the other baselines is a threat of its own: DE is untuned (its low feasibility, \NFeasDE\%, may reflect penalty-dominated search, as for the old PSO setting), MS-SLSQP uses finite differences on a discontinuous objective, and the old PSO setting was run without velocity clamping, which largely restores its feasibility; its effect was measured with one boundary rule and one penalty scaling (Section~\ref{sec:psosetting}). The comparison of SSA and LX-SSA confounds the Laplace step with the doubled cost per iteration, so it tests LX-SSA as published, not the step.
% CHECK (manual, see 05_setup): 2 of 24 Mann-Whitney tests reject, 1.2 expected; CHECK-DIAG [D01]/[D02]: instrumented runs reproduce the stored runs; CHECK-FINAL [C61]: controls switch at 3,015, swarms at 3,030; generated \NFeasDE (R1 editor: "may reflect")
% CHECK-CSWEEP [S10]/[S11]: vmax \NSFeasVMax % feasible ("largely restores"), still worse than constriction; CHECK-FINAL [C17]: LX-SSA as published worse

\emph{Construct validity} (do the measures capture the intended constructs?). At 6,030 evaluations, the ``VNS'' phase is essentially one truncated compass-search descent, so the component analysis measures the value of a start for this local search, not for VNS in general; conclusions about neighborhoods would need larger budgets. The objective is the gross (wake-only) energy of a discretized rose under a top-hat Jensen wake tested at the hub center, with a constant $C_T$ and the benchmark's power curve; its values are benchmark values, and a better objective is not necessarily a better farm. Final quality at a fixed budget is our only performance measure; anytime behavior is shown in the convergence curves but not tested. The feasibility-first ranking limits but does not remove survivorship bias, as the qualification threshold admits methods with up to half of their runs infeasible.
% CHECK-DIAG [D08]-[D10]: truncated compass search; CHECK-FINAL [C21]: cubic curve lowers the power, ranking intact; CHECK-DIR [F03]: coarse bins exploited

\emph{External validity} (do the results generalize?). All results are for discretized roses (15\textdegree{} benchmark bins, 5\textdegree{} on Horns Rev~1), which optimized layouts exploit, circular farms (parallelogram for Horns Rev~1), a $4D$ spacing and a prescribed $N$ of at most 15 on the benchmark and 36 on IEA37, with random starts and 6,030 evaluations unless stated. The rankings are those of a small-budget regime: on the six largest cases at 120,030 evaluations, the leading methods lie within \NBudCloseGapOneTwentyK~pp, and from feasible starts VNS alone ranks first. The real site is a single 16-turbine block; the feasibility advantage of PSO-VNS rests on it. With 1\textdegree{} bins no optimized Horns Rev~1 layout exceeds the installed block; our Horns Rev~1 model takes $C_T$ at the free-stream speed (wake loss \NFPyWakeLossDiffPP~pp relative to PyWake's NOJ model at identical bins). On IEA37, the best PSO-VNS layouts at 30,030 evaluations have AEP gaps of \NFIEAGapSixteenThirtyKProj\% (16 turbines) and \NFIEAGapThirtySixThirtyKProj\% (36 turbines) to the best feasible published layouts, radially projected onto the boundary, and the 64-turbine scenario was not run. The pool lacks CMA-ES~\cite{HansenOstermeier2001} and L-SHADE, the standard references for continuous black-box optimization, so ``best of eight'' means best among the methods compared. Turbulence and loads (critical at $4D$), wake steering, blockage, terrain, cables, electrical losses and costs are outside our scope; all AEP values are gross.
% CHECK-DIR [F03]-[F07]: coarse bins exploited; [F21]: 0 of 901 feasible runs above the installed block at 1 deg; [F32]/[F34]: PyWake 661.39 vs ours 666.75 GWh/yr, wake loss -0.72 pp, cause C_T at free-stream speed; [F44]: projected IEA37 gaps -1.93 / -6.23 % (30,030)
% CHECK-FINAL [C45] (budget study; NBudCloseGapOneTwentyK generated); [C28]: feasible starts, VNS alone first; CHECK-EXTRA [X52]: Horns Rev feasibility (one site, random starts)

\emph{Statistical conclusion validity} (are the inferences sound?). The equivalence margin of $\pm\NEqMargin$~pp was stated after the primary analysis, just above the case-level minimal margin; it is a convention, and the minimal margins (\NXEqSeedPSOVNSvsPSOMin{} to \NXEqClustPSOVNSvsPSOMin~pp) are reported so that readers can apply their own. The 68 cases form only \NXClusters{} (data set, radius) clusters with nested $N$: an exact sign test over the clusters cannot fall below $p=\NXClMinP$, the cluster-robust intervals have 5 degrees of freedom, and the p-values describe this benchmark rather than a population of farms. The $N\ge10$ split, the Data Set~II subgroup and the density pattern were stated after the primary analysis and are exploratory; with one Holm family over all \NXMultN{} case-mean tests, the Data Set~II subgroup remains significant. The budget and feasible-start rankings use six cases and are not tested, and the Horns Rev~1 feasibility result rests on one site with 30 seeds.
% CHECK-FINAL [C49] + CHECK-EXTRA [X41]: post hoc margin, minimal margins; [X10]: six clusters; [X11]: exact p = \NXClMinP; [X42]: CR2 5 d.f.; [X47]: density pattern (exploratory); [X18]/[X19]: all-family Holm, DS II subgroup survives; CHECK-FINAL [C29]: rank differences not tested; CHECK-EXTRA [X52]

\section{Conclusions and Future Work}
\label{sec:conclusion}
On the discretized Kusiak--Song benchmark with random starts and 6,030 evaluations, three findings emerge. First, the baseline configuration changes the conclusions: the old PSO setting, which lies beyond the order-2 boundary and rarely samples feasible layouts late in the run, puts PSO behind salp-swarm hybrids, whereas the constriction setting puts it ahead of every salp-swarm method at this budget. Second, salp-swarm first phases add nothing beyond random sampling in the farm (a gain above the margin is ruled out across seeds and cases; LX-SSA as published is worse), whereas a PSO phase gives the second phase, essentially one compass-search descent, a clearly better start. Third, PSO-VNS has the best average rank with 15\textdegree{}, 5\textdegree{} and 1\textdegree{} direction bins, also under a Gaussian wake at 1\textdegree{}, but its average advantage over a well-configured PSO is small: within $\pm\NEqMargin$~pp under the benchmark rose (across seeds and cases, not farm clusters) and \NFPairMeanOne~pp under a 1\textdegree{} rose, with case-level differences in both directions that follow a data set $\times$ density pattern (exploratory). PSO-VNS reaches feasibility more reliably on Horns Rev~1 (\NHRFeasPSOVNS{} vs.\ \NHRFeasPSO{} runs, McNemar $p=\NXHRMcNemarP$); it is a consistent reference method, not a superior optimizer.
% CHECK-FINAL [C12]/[C27]: old setting puts PSO behind salp-swarm hybrids, constriction ahead of every salp-swarm method (6,030); CHECK-DIAG [D05]: few feasible particle evaluations late in the run; CHECK-CSWEEP [S04]
% CHECK-FINAL [C53]/[C54]: no SSA gain over disc sampling = "in the farm"; [C55]: LX-SSA-VNS worse; [C56]/[C14]/[C44]: PSO phase beats both controls; CHECK-EXTRA [X53]: gain > margin ruled out across seeds and cases (one-sided); CHECK-DIAG [D08]: essentially one descent
% CHECK-DIR [F06]/[F09] + CHECK-FINAL [C21]/[C22]: PSO-VNS first at 15/5/1 deg (Jensen), cubic curves, Gaussian 1 deg; PSO first only under Gaussian 15 deg; CHECK-EXTRA [X39]: seed/case equivalent, cluster not; CHECK-DIR [F10]: -0.043 pp at 1 deg; CHECK-EXTRA [X43]: 19/68 beyond margin, both directions; [X47]: data set x density pattern (exploratory); [X52]: Horns Rev 30/30 vs 19/30, McNemar p 9.8e-4; CHECK-FINAL [C08]/[C20]

The broader lesson is that the conclusions of a metaheuristic comparison are only as reliable as its weakest baseline and its weakest control. Five directions follow from the threats above. \emph{Expensive wake models:} the benchmark's analytical wake takes milliseconds per layout; with engineering or CFD-based wake models, a surrogate-assisted evaluation, for example a graph neural network that maps a layout and a wind state to turbine powers, or the data-driven and Kriging surrogates used in~\cite{Yang2023,Wang2024,Li2025}, would make equal-budget comparisons at realistic fidelity affordable, provided the surrogate error is reported alongside the method differences. \emph{Larger farms:} the IEA37 64-turbine scenario and full farms such as the 80-turbine Horns Rev~1, where a compass-search sweep costs $2n$ evaluations and the budget per coordinate shrinks. \emph{Exact-gradient baselines:} gradient-based methods on a smooth wake model with exact gradients, as in the comparison of~\cite{Thomas2023}, and pseudo-gradients~\cite{Quaeghebeur2021}, in place of finite differences on a discontinuous objective. \emph{Stronger controls and baselines:} CMA-ES, L-SHADE and tuned DE; a cost-matched ablation of the Laplace step; and random-sampling controls at larger budgets and with feasible starts. \emph{Finer roses during optimization}, not only in the re-evaluation, together with economic objectives and loads, so that benchmark gains can be related to design value.
